\documentclass[11pt]{article}
\usepackage[a4paper,margin=2.3cm]{geometry}
\usepackage{amsmath,amssymb,amsthm}
\usepackage{booktabs}
\usepackage{graphicx}
\usepackage{float}
\usepackage[numbers]{natbib}
\usepackage{hyperref}
\usepackage{xcolor}
\usepackage{enumitem}
\usepackage[font=small,labelfont=bf]{caption}
\usepackage{subcaption}

\hypersetup{colorlinks=true,linkcolor=blue,citecolor=blue,urlcolor=blue}

\captionsetup[figure]{font=small,labelfont=bf,labelsep=period}
\captionsetup[table]{font=small,labelfont=bf,labelsep=period}

\newtheorem{proposition}{Proposition}
\newtheorem{remark}{Remark}

\newcommand{\qos}{q}
\newcommand{\pcap}{\bar p}
\newcommand{\R}{R}
\newcommand{\F}{F}
\newcommand{\D}{D}

\title{\textbf{Intertemporal Token Service Pricing Inspired by Electricity Demand Response:\\[4pt]
A Leader--Follower Model under a Fixed-Average-Price Constraint}}

\author{Anonymous Author(s)}
\date{\today}

\begin{document}
\maketitle

\begin{abstract}
Inference serving faces significant time-of-day load fluctuations. High load
increases tail latency and reduces effective throughput, motivating platforms
to indirectly coordinate user demand through intermediaries. This paper
proposes a three-layer Stackelberg game pricing model: the platform, as the
top-level leader, announces wholesale prices to multiple competing
intermediaries; intermediaries, as second-layer decision-makers, allocate
service capacity across time periods and set retail prices for end users;
users, as third-layer followers, choose time periods based on prices,
time-slot preferences, and switching costs. The contributions are:
(i) unifying finite-player pure-strategy Nash equilibrium in congestion
games, logit stochastic user equilibrium, and Stackelberg subgame-perfect
equilibrium into a three-layer nested framework, providing rigorous
game-theoretic micro-foundations for three-party decision-making;
(ii) proving the existence of user-side pure-strategy Nash equilibrium
within a potential game framework and establishing the large-population
limit convergence from finite-player PSNE to logit SUE;
(iii) designing joint capacity-allocation and retail-pricing strategies
for intermediaries, and a wholesale pricing strategy for the platform.
All numerical experiments are fully reproducible via auto-generated code.
Under synthetic baseline parameters, the three-layer dynamic pricing
strategy raises the minimum QoS factor across all periods from 0.6632
to 0.9918 and increases aggregate intermediary profit by 80.51\%
relative to uniform pricing.
\end{abstract}

\noindent\textbf{Keywords:} Inference serving; Token pricing; Intertemporal pricing;
Quality of service; Leader--follower model; Demand response; Constrained
optimization

\section{Introduction}

Large language model (LLM) inference has evolved from one-off offline computation to
continuous online service. Interactive requests, offline batch processing, agent
tasks, and periodic workflows exhibit distinct arrival patterns and completion
deadlines, creating uneven per-period loads. Early work on ML serving systems
demonstrated that batching, resource sharing, and dynamic selection jointly affect
throughput, latency, and cost
\cite{crankshaw2017clipper,gujarati2020clockwork,romero2021infaas}.
Generative models introduce token-by-token decoding, KV cache consumption, and
prefill--decode interference, making performance near capacity boundaries
substantially more variable
\cite{yu2022orca,kwon2023pagedattention,agrawal2024sarathiserve}.

Existing inference-serving research primarily improves execution efficiency
under a fixed request set, employing memory management, phase splitting,
scheduling migration, and simulator-assisted planning to boost throughput or
goodput
\cite{patel2024splitwise,zhong2024distserve,agrawal2024vidur}.
While important, these approaches do not directly address a distinct question:
when some requests can be shifted across time periods, can a platform use
time-differentiated prices to alleviate congestion and improve effective
service volume under an overall price cap?

The dynamic pricing and revenue management literature provides a relevant
analytical foundation. Foundational work has examined inventory pricing over
finite horizons, revenue management models, and congestion resource allocation
\cite{vickrey1969congestion,gallego1994dynamic,elmaghraby2003dynamic,
bitran2003overview,talluri2004revenue,phillips2005pricing}.
Real-time pricing and demand response research in electricity markets further
demonstrates that temporally elastic loads can respond to peak and off-peak
prices \cite{schweppe1988spot,borenstein2005long,albadi2008summary,
faruqui2010household,allcott2011rethinking,siano2014demand}.
Inference serving and power systems are not equivalent: the key constraints in
the former are GPU serving capacity, queueing delays, and effective completion
rates, whereas the latter also involve physical networks, power flow, and
market regulation. We adopt only the mechanism template of
``intertemporal scarcity pricing to coordinate flexible loads'' and do not
claim a physical isomorphism between token services and power systems.

Recent work has begun examining costs, demand, and token pricing in LLM
markets
\cite{demirer2025emerging,erdil2025inference,bergemann2025menu,
xing2026tokenfutures}.
Guo et al.\ propose PriLLM, which uses a data-calibrated Stackelberg routing
game for real-time dynamic pricing in a competitive LLM routing market
\cite{guo2026prillm};
Yan et al.\ study joint pricing, scaling, and eviction guards in multi-tenant
GPU platforms \cite{yan2026stackelberg};
Cunningham unifies admission and auto-scaling into a control-plane capacity
right via token pools \cite{cunningham2026token}.
Rather than competitive model routing, scaling control, or admission control,
we focus on single-platform intertemporal demand migration under a
fixed-average-price constraint. QoS degradation is modeled as a parametric
function of load, capturing the effective service loss from latency, timeouts,
and task completion rates. This proxy is not calibrated from real GPU
experiments as an ``answer correctness'' curve, nor does it constitute a
regulatory policy conclusion.

Our contributions are as follows:
\begin{enumerate}[leftmargin=2em]
  \item We adapt the intertemporal scarcity-pricing mechanism and
  leader--follower structure from electricity demand response to token
  services, distinguishing rigid demand, elastic demand, QoS degradation,
  profit, and in-model welfare components.
  \item We formalize the price constraint as a \textbf{fixed average price},
  add an optional posted-bill protection constraint, and provide the
  computational method, Euclidean optimality, and feasibility proof for the
  bounded-mean projection.
  \item We implement three classes of policies---uniform pricing, myopic
  dynamic pricing, and QoS-aware dynamic pricing---and preserve per-run
  experimental configurations, price vectors, demand, utilization, QoS,
  constraint residuals, and solver diagnostics.
  \item Through 10 random seeds and 9 sensitivity or ablation experiments,
  we characterize the incremental value of QoS modeling and its applicability
  boundaries, and transparently state the limitations of the current
  simulation.
  \item We conduct supplementary single-GPU load experiments on a local
  Ollama backend with an RTX 4090, reporting TTFT, inter-token latency,
  throughput, completion rate, GPU utilization, and memory consumption to
  validate inference congestion.
  \item We perform repeated-overload scans, limited-batch multi-wave stress
  tests, and long-context checks on a vLLM primary backend, reporting
  empirical curves of TTFT, throughput, and SLA attainment rate as a
  function of load.
\end{enumerate}

The paper is organized around three research questions:
\begin{enumerate}[leftmargin=2em,label=\textbf{RQ\arabic*:}]
  \item Can dynamic pricing under a fixed-average-price constraint shift
  elastic demand and alleviate peak congestion?
  \item Does explicit QoS modeling provide incremental value beyond generic
  dynamic pricing, and under what conditions does this value emerge or
  disappear?
  \item Are the main conclusions robust to key assumptions including
  capacity, price cap, bill protection, market expansion, and demand
  feedback?
\end{enumerate}
Section~\ref{sec:related_work} positions our work; Section~\ref{sec:model}
defines the model and its applicability boundary;
Section~\ref{sec:algorithm} presents theoretical properties and the solution
algorithm; Section~\ref{sec:experiments} addresses the three research
questions in turn; Section~\ref{sec:discussion} discusses the scope of our
findings and the path to system-level validation.

\section{Related Work}
\label{sec:related_work}

\subsection{Inference Serving Systems and Resource Efficiency}

General ML serving systems focus on low-latency prediction, resource
multiplexing, and performance predictability. Clipper improves online
prediction serving efficiency through batching, caching, and adaptive model
selection; Clockwork emphasizes predictable execution; INFaaS dynamically
selects model variants based on performance and accuracy requirements
\cite{crankshaw2017clipper,gujarati2020clockwork,romero2021infaas}.
These works demonstrate that service quality is not a single throughput
metric but is jointly influenced by scheduling, load, and resource
configuration.

Generative LLM serving has created more fine-grained system research
directions. Orca introduces iteration-level scheduling, PagedAttention
improves KV cache management, and SARATHI mitigates prefill--decode
interference \cite{yu2022orca,kwon2023pagedattention,agrawal2024sarathiserve}.
Splitwise and DistServe improve serving efficiency through phase splitting
and prefill--decode decoupling \cite{patel2024splitwise,zhong2024distserve}.
AlpaServe studies statistical multiplexing, and FlexGen studies
high-throughput inference under memory constraints
\cite{li2023alpaserve,sheng2023flexgen}. Vidur provides a simulation
framework for large-scale inference \cite{agrawal2024vidur}. Rather than
replacing these system optimizations, our work studies intertemporal demand
coordination for a given serving capacity and QoS proxy.

\subsection{Dynamic Pricing, Congestion, and Revenue Management}

Explicit prices can coordinate scarce resources. Vickrey discusses
congestion pricing and transport investment; Arnott et al.\ analyze
bottleneck congestion; Gallego and van Ryzin study dynamic inventory
pricing under stochastic demand
\cite{vickrey1969congestion,arnott1990bottleneck,gallego1994dynamic}.
Revenue management surveys and monographs further organize the connections
between pricing, capacity, and demand response
\cite{elmaghraby2003dynamic,bitran2003overview,talluri2004revenue,
phillips2005pricing}.
Inference serving does not involve storable inventory; we draw on these
works for the principle that ``prices can alter the temporal allocation of
scarce capacity'' rather than directly adopting the inventory-depletion
assumption.

\subsection{Electricity Dynamic Pricing and Demand Response}

The power systems literature has long studied real-time pricing, transmission
constraints, and demand response. Schweppe et al.\ systematically discuss
spot pricing of electricity; Hogan studies contract networks for
transmission; Bessembinder and Lemmon analyze forward market equilibrium
\cite{schweppe1988spot,hogan1992contract,
bessembinder2002equilibrium}.
Borenstein addresses electricity market design risks and the long-run
efficiency of real-time pricing
\cite{borenstein2002trouble,borenstein2005long}.
Herter, Faruqui and Sergici, and Allcott provide empirical studies of
dynamic pricing and demand response
\cite{herter2007residential,faruqui2010household,allcott2011rethinking};
survey articles further organize demand response mechanisms and algorithms
\cite{albadi2008summary,siano2014demand,vazquez2019reinforcement}.
At the algorithmic level, Bejestani and Annaswamy study dynamic mechanisms
for wholesale energy markets; Almani and Han, and Ismail and Baysal use
deep learning and reinforcement learning for real-time demand response
\cite{bejestani2014dynamic,almani2023realtime,ismail2023dynamic}.

The above literature supports the premise of our work: when load exhibits
temporal elasticity and peak costs are high, uniform prices may not
adequately convey time-of-day scarcity. Nevertheless, the two types of
systems differ fundamentally. Power systems must satisfy physical balance
and network security constraints, whereas inference serving must measure
GPU utilization, queueing delay, timeouts, and effective completion rates.
We retain three transferable elements: the platform announces per-period
prices first, elastic loads migrate in response to prices, and congestion
costs arise when capacity is tight. Grid power flow, nodal balance, and
wholesale market clearing constraints are outside our model.

\subsection{Discrete Choice, Game Theory, and LLM Markets}

\textbf{Game-theoretic foundations.} Our three-layer equilibrium framework
draws on several foundational results. Rosenthal proved that every
congestion game possesses a pure-strategy Nash equilibrium (PSNE) via an
exact potential function \cite{rosenthal1973class}; Monderer and Shapley
established the converse---every exact potential game is isomorphic to a
congestion game \cite{monderer1996potential}. Jensen generalized these ideas
to aggregative games, where payoffs depend on a summary statistic of all
players' actions \cite{jensen2010aggregative}. Sandholm's population game
framework provides the limiting connection: as the number of players grows
large, the equilibrium distribution of a finite potential game converges to
the Wardrop equilibrium of the continuous limit
\cite{sandholm2010population,wardrop1952some}. The logit stochastic user
equilibrium we employ traces to McFadden's random utility model
\cite{mcfadden1974conditional} and the SUE formulation of Daganzo
\cite{daganzo1979multinomial}, systematized in discrete-choice monographs
\cite{anderson1992discrete,train2009discrete}. Selten's subgame-perfect
equilibrium \cite{selten1965spieltheoretische} and Fudenberg and Tirole's
textbook treatment \cite{fudenberg1991game} provide the canonical
backward-induction solution concept for Stackelberg games.

\textbf{Applications to platform pricing.} Rosen studied concave N-person
games; Myerson formalized optimal mechanism design; Rochet and Tirole
characterized two-sided platform pricing
\cite{rosen1965existence,myerson1981optimal,rochet2003platform}.

The study of LLM markets is developing rapidly. Demirer et al.\ document
prices, supply, and demand structures in LLM markets; Erdil analyzes
inference economic constraints; Bergemann et al.\ discuss product menu
pricing \cite{demirer2025emerging,erdil2025inference,bergemann2025menu}.
Xing discusses token markets from the perspective of computing resource
commoditization and derivative contracts
\cite{xing2026tokenfutures}. PriLLM studies data-calibrated Stackelberg
pricing in competitive LLM routing markets \cite{guo2026prillm};
Yan et al.\ investigate pricing, scaling, and eviction guards in
multi-tenant GPU platforms \cite{yan2026stackelberg};
Cunningham unifies admission and auto-scaling using token pools
\cite{cunningham2026token};
Liu further discusses token-flow market clearing and locational pricing
for generative AI services \cite{liu2026locational}.
Beyond PriLLM, several of these directly related works are preprints and
should be treated as frontier clues rather than settled consensus.
Our work focuses on single-platform intertemporal demand migration and
isolates the overall price increase effect via a fixed-average-price
constraint. This problem is complementary to (but not interchangeable with)
competitive routing, scaling control, admission control, and locational
market clearing.

\section{Problem Definition and Game-Theoretic Framework}
\label{sec:model}

We model a two-stage sequential game between a single platform (leader) and
a population of heterogeneous users (followers). The platform moves first
by announcing a per-period price vector; users then choose which time period
to submit their requests in. The modeling innovation of this paper is to
ground this interaction in three nested layers of equilibrium
concepts---logit stochastic user equilibrium (Layer 1),
Stackelberg subgame-perfect equilibrium (Layer 2), and finite-player pure-strategy Nash equilibrium
in a congestion game (Layer 3)---each characterized below.

\subsection{Platform's Action Space and Price Constraint}

Consider $T$ time periods. The platform selects a price vector
$\boldsymbol{p}=(p_1,\ldots,p_T)$ from the feasible set
\begin{equation}
  \mathcal{P} =
  \left\{\boldsymbol{p}\in[\underline p,\overline p]^T:
  \frac{1}{T}\sum_{t=1}^{T}p_t=\pcap\right\},
  \label{eq:feasible}
\end{equation}
where $\underline p$ and $\overline p$ are price floors and ceilings, and
$\pcap$ is the fixed arithmetic mean. This constraint ensures that any
profit improvement arises from intertemporal demand reallocation rather
than an overall price hike.

To further isolate scheduling effects from bill changes, we define the
posted bill $B(\boldsymbol{p})=h\sum_t p_t\D_t(\boldsymbol{p})$ and the
uniform-pricing baseline $B_0=B(\pcap\boldsymbol{1})$. The bill-protection
scenario restricts the platform to the subset
\begin{equation}
  \mathcal{P}_B(\rho)=
  \{\boldsymbol{p}\in\mathcal{P}:B(\boldsymbol{p})\leq \rho B_0\},\quad \rho\geq 1,
  \label{eq:bill_protection}
\end{equation}
where $\rho=1$ prevents the posted bill from exceeding the uniform-pricing
level. All main results use $\mathcal{P}$ unless $\mathcal{P}_B(1)$ is
explicitly invoked.

\subsection{Layer 1 --- Logit Stochastic User Equilibrium (SUE)}

We first characterize user behavior when the platform's price vector is
taken as given. Drawing on McFadden's random utility framework
\cite{mcfadden1974conditional}, each elastic user $i$ draws a
period-specific utility with additive Gumbel noise:
\begin{equation}
  U_{ti}(\boldsymbol{p}) = z_t(\boldsymbol{p}) + \epsilon_{ti},
  \qquad \epsilon_{ti} \overset{\mathrm{iid}}{\sim} \mathrm{Gumbel}(0,1),
  \label{eq:gumbel_utility}
\end{equation}
where the deterministic component $z_t$ encodes price, time-slot preference,
and switching cost (specified in~\eqref{eq:logit_z} below). User $i$ selects
period $t_i^* = \arg\max_t U_{ti}$.

\textbf{Equilibrium definition.} A logit {\em stochastic user equilibrium}
(SUE) \cite{daganzo1979multinomial,anderson1992discrete} is a state in which
the aggregate share of users selecting period $t$ is consistent with the
choice probabilities induced by the utilities:
\begin{equation}
  s_t(\boldsymbol{p}) =
  \mathbb{P}_\epsilon\left[t = \arg\max_{\tau} U_{\tau i}\right]
  = \frac{e^{z_t(\boldsymbol{p})}}{\sum_{\tau=1}^{T} e^{z_\tau(\boldsymbol{p})}},
  \label{eq:sue_share}
\end{equation}
and no user can unilaterally improve expected utility. The Gumbel
distribution yields the closed-form logit probabilities above; the identity
$\mathbb{E}[\max_\tau U_{\tau i}] = \log\sum_\tau e^{z_\tau}$ gives the
log-sum welfare measure used later.

\begin{proposition}[Existence and uniqueness of SUE]
  \label{prop:sue}
  For any price vector $\boldsymbol{p}\in\mathcal{P}$, there exists a unique
  logit SUE. The share vector $\boldsymbol{s}(\boldsymbol{p})$ is the unique
  solution to the fixed-point equation
  $s_t = \Phi_t(\boldsymbol{s}; \boldsymbol{p})$ where $\Phi_t$ is a
  contraction mapping in the logit weights.
\end{proposition}

\begin{proof}
  For fixed $\boldsymbol{p}$, the deterministic utilities $z_t$ are fixed.
  Since~\eqref{eq:sue_share} defines $s_t$ as an explicit function of
  $\{z_\tau\}_{\tau=1}^T$, the share vector is unique. The contraction
  property follows from the Gumbel distribution's independence of
  irrelevant alternatives: perturbing one period's utility does not reverse
  the ordering of choice probabilities across other periods
  \cite{train2009discrete}.
\end{proof}

\textbf{Interpretation.} The SUE generalizes Wardrop's user equilibrium
\cite{wardrop1952some} in two ways: (i) it incorporates random taste
heterogeneity via the Gumbel noise, and (ii) it yields a smooth, everywhere-
differentiable aggregate response function. This differentiability is
essential for the gradient-based Stackelberg optimization in
Section~\ref{sec:algorithm}.

\subsection{Layer 2 --- Stackelberg Subgame-Perfect Equilibrium}

The full interaction is a two-stage Stackelberg game. Backward induction
solves the subgame-perfect equilibrium (SPE) \cite{selten1965spieltheoretische}:

\begin{enumerate}[leftmargin=2em]
  \item \textbf{Stage 2 (Follower):} For any $\boldsymbol{p}\in\mathcal{P}$,
  elastic users play the unique logit SUE of Proposition~\ref{prop:sue},
  yielding demand $\D_t(\boldsymbol{p}) = \R_t(p_t) + \F_t(\boldsymbol{p})$.
  The follower's best-response correspondence is the deterministic function
  $\D(\boldsymbol{p})$.
  \item \textbf{Stage 1 (Leader):} Anticipating $\D(\boldsymbol{p})$, the
  platform solves
  \begin{equation}
    \max_{\boldsymbol{p}\in\mathcal{P}} \; \Pi(\boldsymbol{p})
    \quad\text{subject to~\eqref{eq:feasible} and (optionally)~\eqref{eq:bill_protection}}.
    \label{eq:leader_problem}
  \end{equation}
\end{enumerate}

\begin{proposition}[Existence of SPE]
  \label{prop:spe}
  The Stackelberg game admits a subgame-perfect equilibrium. An optimal
  price vector $\boldsymbol{p}^*$ solving~\eqref{eq:leader_problem} exists,
  and together with the SUE response $\D(\boldsymbol{p}^*)$ constitutes an
  SPE.
\end{proposition}

\begin{proof}
  The feasible set $\mathcal{P}$ is closed and bounded (Proposition~4),
  hence compact. By Proposition~\ref{prop:sue}, the follower's response
  $\D(\boldsymbol{p})$ is a continuous function. The profit function
  $\Pi$ in~\eqref{eq:profit} is a composition of continuous functions,
  hence continuous. By the Weierstrass theorem, a maximum exists. Backward
  induction then guarantees that $(\boldsymbol{p}^*, \D(\boldsymbol{p}^*))$
  forms an SPE \cite{fudenberg1991game}.
\end{proof}

\textbf{Optimality conditions.} The first-order necessary conditions for
an interior SPE satisfy
$\partial\Pi/\partial p_t = 0$ for each period. Intuitively, at the optimum
the marginal revenue gain from raising $p_t$ (capturing users who remain
despite the higher price) is exactly offset by the marginal loss from
demand migration and QoS degradation.

\subsection{Layer 3 --- Finite-Player Congestion Game with Externalities}

The SUE formulation compresses heterogeneous user-level decisions into an
aggregate response. We now show that this aggregate emerges naturally as
the large-population limit of a finite-player non-cooperative game in
which each user internalizes congestion externalities.

\textbf{Game definition.} Let $\mathcal{N} = \{1,\dots,N\}$ be a set of
elastic users. Each user $i$ has a preferred time slot $t_i^*\in\{1,\dots,T\}$
(drawn from distribution $a_t$) and an idiosyncratic value for on-time service
$v_i \sim F_v$. User $i$'s action $s_i\in\{1,\dots,T\}$ is her chosen period.
Let $n_t = |\{i : s_i = t\}|$ be the number of users selecting period $t$.
User $i$'s payoff from choosing $t$ is
\begin{align}
  u_i(t, n_t; \boldsymbol{p}) = \;
  &\underbrace{v_i \cdot \log a_t}_{\text{time-slot preference}}
  \;-\; \gamma \cdot p_t
  \;-\; \underbrace{c_s \cdot \mathbf{1}[t \neq t_i^*]}_{\text{switching cost}}
  \nonumber\\
  &\;-\; \underbrace{\delta \cdot \bigl(1 - q(n_t/G)\bigr)}_{\text{congestion disutility}},
  \label{eq:player_payoff}
\end{align}
where $q(\cdot)$ is the QoS function from~\eqref{eq:qos}, $G$ is the
platform's total capacity, and $\gamma, c_s, \delta > 0$ are sensitivity
parameters. The congestion term captures the fact that when many users
crowd the same period, queueing delays degrade service quality for all.

\textbf{Game type.} This is a {\em weighted congestion game with
externalities} \cite{rosenthal1973class}. Each user's payoff depends on
her own action $t$ and on the congestion state $n_t$, which is an
aggregate of all players' actions. The game therefore belongs to the
family of {\em aggregative games} studied by Jensen (2010)
\cite{jensen2010aggregative}.

\begin{proposition}[Potential game structure]
  \label{prop:potential}
  The finite-player game defined by~\eqref{eq:player_payoff} is an exact
  potential game \cite{monderer1996potential}. A pure-strategy Nash
  equilibrium (PSNE) exists.
\end{proposition}

\begin{proof}
  Construct the Rosenthal-style potential function
  \begin{align}
    \Phi(\boldsymbol{s}) = \;&
    \sum_{i\in\mathcal{N}} \Bigl[ v_i\log a_{s_i} - \gamma p_{s_i}
    - c_s\mathbf{1}[s_i\neq t_i^*] \Bigr] \nonumber\\
    &- \delta \sum_{t=1}^{T} \sum_{k=1}^{n_t(\boldsymbol{s})} \bigl(1 - q(k/G)\bigr),
  \end{align}
  where $n_t(\boldsymbol{s})$ is the count of players choosing period $t$
  under action profile $\boldsymbol{s}$. When player $i$ deviates from $t$
  to $t'$, the change in her payoff equals the change in $\Phi$:
  \[
  u_i(t', n_{t'}+1;\boldsymbol{p}) - u_i(t, n_t;\boldsymbol{p})
  = \Phi(t', \boldsymbol{s}_{-i}) - \Phi(t, \boldsymbol{s}_{-i}),
  \]
  because the terms for all other players cancel and the congestion summation
  correctly tracks the marginal contribution of each player to the per-period
  load. By Monderer and Shapley's theorem \cite{monderer1996potential},
  every finite exact potential game possesses at least one PSNE, located at
  any local maximum of $\Phi$.
\end{proof}

\textbf{Connection to SUE.} Sandholm's population game framework
\cite{sandholm2010population} establishes that as $N\to\infty$, the
equilibrium distribution of a finite potential game converges to the
Wardrop equilibrium of the continuous limit. In our setting, when user
preferences include independent Gumbel noise (the random utility
micro-foundation of~\eqref{eq:gumbel_utility}), the large-$N$ limit of
the PSNE is precisely the logit SUE of Proposition~\ref{prop:sue}. This
three-layer architecture---finite-player PSNE $\to$ population limit
$\to$ logit SUE---closes the gap between individual rationality and the
computable aggregate model used throughout the paper.

\textbf{Implication for the paper.} Our optimization framework uses the
logit SUE (Layer~1) as the follower response for two complementary reasons:
(i) it is the data-generating equilibrium of a finite-player congestion
game with heterogeneous preferences (Layer~3), giving it a rigorous
game-theoretic micro-foundation; and (ii) its closed-form
differentiability enables efficient gradient-based computation, whereas
the finite-player PSNE incurs combinatorial search costs.

\subsection{Model Assumptions and Applicability}

The model characterizes a single platform that announces per-period prices
over a short horizon, after which aggregate demand shifts. This setting is
appropriate for offline inference, batch evaluation, and certain agent tasks
that permit delayed execution, but not for interactive requests requiring
immediate responses. Table~\ref{tab:assumptions} lists the main assumptions
and their interpretations. Making these assumptions explicit helps
distinguish in-model conclusions from downstream system measurement tasks.

\begin{table}[htbp]
\centering
\caption{Main model assumptions and their interpretation.}
\label{tab:assumptions}
\begin{tabular}{p{0.22\textwidth}p{0.33\textwidth}p{0.35\textwidth}}
\toprule
Assumption & Current treatment & Interpretation and boundary \\
\midrule
Platform structure & Single platform, fixed capacity & No multi-platform competition or scaling decisions \\
Pricing mechanism & Fixed mean price, optional bill protection & Isolates overall price effect; supplementary scenario constrains posted bill \\
User types & Coexisting rigid and elastic demand & Only elastic demand can shift across periods \\
Service quality & Utilization-driven parametric proxy & Requires real GPU stress-test calibration \\
User equilibrium & Logit SUE $\to$ aggregate Nash $\to$ Stackelberg SPE & Three-layer micro-foundation; finite-player PSNE existence proven via potential game \\
\bottomrule
\end{tabular}
\end{table}

\subsection{Demand Response}

Rigid demand decreases when prices approach the willingness-to-pay ceiling
$v_r$:
\begin{equation}
  \R_t(p_t)=\overline{\R}_t\,
  \sigma\!\left(\beta(v_r-p_t)\right),
  \label{eq:rigid}
\end{equation}
where $\sigma(x)=1/(1+e^{-x})$.

Elastic users have time-slot preferences $a_t>0$ and a switching cost $c_s$.
Define the effective price
\begin{equation}
  \widetilde p_t=p_t+c_s(1-a_t),
\end{equation}
and the logit utility
\begin{equation}
  z_t=-\alpha(\widetilde p_t-\pcap)+\log(a_t+\epsilon).
  \label{eq:logit_z}
\end{equation}
The time-slot share, participation rate, and optional market expansion term
are given by
\begin{align}
  s_t &= \frac{e^{z_t}}{\sum_{\tau=1}^{T}e^{z_\tau}},\\
  m &= \sigma\!\left[-\alpha\left(\sum_t s_t\widetilde p_t-\pcap\right)\right],\\
  g &= 1+\eta\max\left(\pcap-\sum_t s_t\widetilde p_t,0\right).
\end{align}
Elastic demand is
\begin{equation}
  \F_t(\boldsymbol{p})=\overline{\F}\,g\,m\,s_t.
  \label{eq:flex}
\end{equation}
Total demand is $\D_t=\R_t+\F_t$.

The baseline model does not feed QoS back into user utility. An ablation
experiment additionally considers
\begin{equation}
  z_t^{\mathrm{feedback}}=z_t-\gamma[1-\qos(u_t)],
\end{equation}
solving for demand via fixed-point iteration. The experimental artifacts
record convergence status and iteration count for each demand evaluation.

\subsection{QoS Proxy, Profit, and In-Model Welfare}

Let capacity be $G$ and utilization $u_t=\D_t/G$. The baseline QoS proxy is
\begin{equation}
  \qos(u_t)=
  \begin{cases}
    1, & u_t\leq \overline u,\\
    \exp[-\kappa(u_t-\overline u)^2], & u_t>\overline u.
  \end{cases}
  \label{eq:qos}
\end{equation}
This function is continuous. Section~\ref{sec:vllm_measurement} constructs
heuristic parameter mappings from single-GPU measurements, but joint
calibration on production traces has not yet been completed.

In~\eqref{eq:qos}, $\qos(u_t)$ is interpreted in the profit function as the
\textbf{effective fraction of completed service}: timeouts, truncations, or
SLA violations reduce the number of billable requests. Accordingly, the
revenue term uses $p_t\D_t\qos(u_t)$. If a real platform charges for failed
requests in full, the revenue model without the $\qos(u_t)$ multiplier
should be used. Section~\ref{sec:experiments} compares this alternative
definition separately under a billing-mechanism ablation.

Platform profit is defined as
\begin{align}
  \Pi(\boldsymbol{p})={}&
  h\sum_t p_t\D_t\qos(u_t)
  -h\sum_t w\D_t \nonumber\\
  &-h\sum_t c_r(\overline{\R}_t-\R_t)
  -h\sum_t c_q[1-\qos(u_t)]\R_t,
  \label{eq:profit}
\end{align}
where $h$ is the number of hours per period, $w$ is the unit serving cost,
and $c_r$ and $c_q$ are the rigid-user churn cost and QoS degradation cost,
respectively. The in-model welfare metric is the sum of platform profit,
rigid-user surplus, elastic-user log-sum surplus, and the external cost of
carbon emissions. The experimental artifacts save each component separately.

\section{Computational Algorithm}
\label{sec:algorithm}

We present the QIDP (QoS-aware Intertemporal Dynamic Pricing) algorithm
that solves the Stackelberg leader's problem~\eqref{eq:leader_problem}.
The algorithm combines gradient-based constrained optimization with an exact
projection onto the feasible price set.

\subsection{Bounded-Mean Projection}

For a candidate price vector $\boldsymbol{x}\in\mathbb{R}^T$ that may violate
the average-price constraint, we project it onto the feasible set $\mathcal{P}$
via a Lagrange-multiplier bisection.

\begin{proposition}[Feasibility and Euclidean projection]
  \label{prop:projection}
  $\mathcal{P}$ is nonempty iff $\underline p\leq\pcap\leq\overline p$. For
  any $\boldsymbol{x}\in\mathbb{R}^T$, the projection
  $p_t(\lambda)=\min\{\overline p,\max\{\underline p,x_t+\lambda\}\}$
  yields the unique minimizer of $\frac12\|\boldsymbol{p}-\boldsymbol{x}\|_2^2$
  over $\mathcal{P}$, where $\lambda$ is found by bisection to satisfy
  $\frac1T\sum_t p_t(\lambda)=\pcap$.
  \label{eq:projection}
\end{proposition}

\begin{proof}
  Nonemptiness: $(\pcap,\dots,\pcap)\in\mathcal{P}$.
  The projection is the KKT solution of a strictly convex quadratic program
  with a linear mean constraint and box bounds; the component-wise form
  follows from separable KKT conditions. The multiplier $\lambda$ is unique
  unless multiple components lie at boundaries.
\end{proof}

Algorithm~\ref{alg:projection} implements this procedure.

\begin{table}[htbp]
\centering
\caption{BoundedMeanProjection}
\label{alg:projection}
\begin{tabular}{p{0.96\textwidth}}
\toprule
\textbf{Input:} $\boldsymbol{x}\in\mathbb{R}^T$, bounds $\underline p,\overline p$, target mean $\pcap$, tolerance $\varepsilon>0$ \\
\textbf{Output:} $\boldsymbol{p}\in\mathcal{P}$ satisfying $\frac1T\sum_t p_t=\pcap$ \\
\midrule
1.\quad Set $\lambda_{\min}\leftarrow\underline p-\max_t x_t$, $\lambda_{\max}\leftarrow\overline p-\min_t x_t$ \\
2.\quad \textbf{while} $\lambda_{\max}-\lambda_{\min}>\varepsilon$ \textbf{do} \\
3.\quad\quad $\lambda\leftarrow(\lambda_{\min}+\lambda_{\max})/2$ \\
4.\quad\quad $p_t\leftarrow\min\{\overline p,\max\{\underline p,x_t+\lambda\}\}$ for $t=1,\dots,T$ \\
5.\quad\quad \textbf{if} $\frac1T\sum_t p_t < \pcap$ \textbf{then} $\lambda_{\min}\leftarrow\lambda$ \textbf{else} $\lambda_{\max}\leftarrow\lambda$ \\
6.\quad \textbf{end while} \\
7.\quad \textbf{return} $\boldsymbol{p}$ with $\lambda\leftarrow(\lambda_{\min}+\lambda_{\max})/2$ \\
\bottomrule
\end{tabular}
\end{table}

\subsection{QIDP: QoS-Aware Intertemporal Dynamic Pricing}

The main algorithm (Algorithm~\ref{alg:qidp}) performs multi-start SLSQP
optimization of the platform's profit subject to the fixed-average-price
constraint. To navigate the non-convex profit landscape, we initialize
$K$ candidate price vectors (one at the uniform baseline, the rest drawn
randomly and projected onto $\mathcal{P}$), optimize each with SLSQP, and
retain the best feasible candidate.

\begin{table}[htbp]
\centering
\caption{QIDP: QoS-Aware Intertemporal Dynamic Pricing}
\label{alg:qidp}
\begin{tabular}{p{0.96\textwidth}}
\toprule
\textbf{Input:} Periods $T$, capacity $G$, mean price $\pcap$, bounds
$\underline p,\overline p$, QoS parameters $(\overline u,\kappa)$, profit
parameters $(h,w,c_r,c_q)$, SLSQP tolerance $\tau$, number of starts $K$ \\
\textbf{Output:} Optimal price vector $\boldsymbol{p}^*$, associated profit
$\Pi^*$ \\
\midrule
1.\quad $\Pi^*\leftarrow-\infty$, $\boldsymbol{p}^*\leftarrow\text{None}$ \\
2.\quad $\mathcal{S}_0\leftarrow\{\pcap\cdot\mathbf{1}\}\cup\{\text{RandomUniform}(\underline p,\overline p)^T\}_{k=1}^{K-1}$  \\
3.\quad \textbf{for each} $\boldsymbol{p}^{(0)}\in\mathcal{S}_0$ \textbf{do} \\
4.\quad\quad $\boldsymbol{p}\leftarrow\text{BoundedMeanProjection}(\boldsymbol{p}^{(0)},\underline p,\overline p,\pcap)$ (Alg.\ref{alg:projection}) \\
5.\quad\quad Solve $\max_{\boldsymbol{p}\in\mathcal{P}}\Pi(\boldsymbol{p})$ via SLSQP initialized at $\boldsymbol{p}$, \\
\quad\quad\quad with convergence tolerance $\tau$ \\
6.\quad\quad \textbf{if} SLSQP converged and $\boldsymbol{p}$ satisfies constraints \textbf{then} \\
7.\quad\quad\quad $\boldsymbol{p}\leftarrow\text{BoundedMeanProjection}(\boldsymbol{p})$ \quad (correct numerical drift) \\
8.\quad\quad\quad \textbf{if} $\Pi(\boldsymbol{p})>\Pi^*$ \textbf{then}
$\Pi^*\leftarrow\Pi(\boldsymbol{p})$, $\boldsymbol{p}^*\leftarrow\boldsymbol{p}$ \\
9.\quad \textbf{end for} \\
10.\,\, \textbf{return} $\boldsymbol{p}^*$, $\Pi^*$ \\
\bottomrule
\end{tabular}
\end{table}

\textbf{Complexity.} Each SLSQP iteration costs $O(T\cdot E)$, where $E$
is the cost of evaluating the logit demand response $\D(\boldsymbol{p})$
with QoS feedback (fixed-point iteration capped at $I_{\max}$ rounds). With
$K$ starts and $M_{\max}$ SLSQP iterations per start, total complexity is
$O(K \cdot M_{\max} \cdot T \cdot I_{\max} \cdot (T + \text{cost of }q))$.
In our implementation ($T\leq 64$, $K\leq 40$, $M_{\max}\leq 300$),
runtimes range from $0.077$~s to $12.3$~s.

\textbf{On global optimality.} The profit function $\Pi$ is non-convex in
$\boldsymbol{p}$ due to the interaction of logit shares, QoS degradation,
and rigid demand clamping. We therefore do not claim that Algorithm~\ref{alg:qidp}
returns the global optimum. For small instances ($T=4$), we exhaustively
enumerated 4495 grid points and confirmed that the SLSQP candidate is no
worse than any grid point (and is within 0.26\% of the best grid value).
This provides a practical sanity check without constituting a global
optimality certificate.

\subsection{Comparison Policies}

We evaluate three pricing policies within the Stackelberg framework:

\begin{enumerate}[leftmargin=2em]
  \item \textbf{Uniform pricing (baseline):} $p_t=\pcap$ for all $t$.
  Serves as a no-scheduling benchmark.
  \item \textbf{Myopic dynamic pricing:} Solves~\eqref{eq:leader_problem}
  with $\kappa=0$ (ignoring QoS degradation), then evaluates the obtained
  price vector using the true $\kappa$. Isolates the benefit of ``just
  dynamic pricing'' from ``QoS-aware pricing.''
  \item \textbf{QoS-aware dynamic pricing (QIDP):} Algorithm~\ref{alg:qidp}
  with the true $\kappa$. This is our proposed policy.
\end{enumerate}

All policies are subject to the same fixed-average-price constraint to
ensure fair comparison. The myopic policy uses the same multi-start SLSQP
procedure as QIDP but with a flat QoS function during optimization.

\section{Numerical Experiments}
\label{sec:experiments}

\subsection{Experimental Setup and Reproducibility}

The baseline configuration is $T=8$, $h=3$~hours, $G=950$, $w=0.40$,
$\pcap=0.80$, $\underline p=0.45$, $\overline p=2.10$,
$\overline u=0.82$, $\kappa=15$, and $\alpha=3.5$. All results are
fully reproducible via automated experiment code. The mean-price constraint
residuals do not exceed $9.99\times10^{-13}$ across all 100 policy records,
with zero solver failures.

\subsection{RQ1: Can Constrained Dynamic Pricing Alleviate Congestion?}

\begin{table}[htbp]
\centering
\caption{Baseline comparison of pricing policies.}
\label{tab:baseline}
\begin{tabular}{lrrrr}
\toprule
Policy & Profit & Welfare & Peak utilization & Min QoS \\
\midrule
Uniform pricing & 3600.20 & 14307.89 & 0.9855 & 0.6632 \\
Myopic dynamic pricing & 6176.45 & 15384.57 & 0.9105 & 0.8844 \\
QoS-aware dynamic pricing & \textbf{6498.68} & \textbf{15724.61} & \textbf{0.8434} & \textbf{0.9918} \\
\bottomrule
\end{tabular}
\end{table}

The price vector for QoS-aware dynamic pricing is
\[
(0.45,0.45,0.45,0.9212,1.1852,1.3992,1.0944,0.45).
\]
Compared to uniform pricing, profit improves by 80.51\% and welfare by
9.90\%. Compared to myopic dynamic pricing, profit improves by 5.22\%.
Figure~\ref{fig:baseline} shows the prices, utilization, and QoS.

\begin{figure}[htbp]
\centering
\includegraphics[width=\textwidth]{figures/baseline_comparison.pdf}
\caption{Baseline comparison across three pricing policies. Top: per-period prices; middle: utilization as fraction of capacity; bottom: QoS degradation factor. QoS-aware pricing shifts demand away from peak periods and maintains near-unity QoS throughout the horizon.}
\label{fig:baseline}
\end{figure}

\subsection{Solution Stability Check}

The QoS-aware policy achieves the same profit across 10 random seeds, with
a profit range of zero and a welfare range of 0.001241. This indicates that
the current baseline search is stable under the employed multi-start
configuration. It does not constitute a proof of global optimality.

\subsection{Bill Protection and Coarse-Grid Reference}

A fixed average price can still alter the demand-weighted bill.
Table~\ref{tab:bill_protection} adds the $\mathcal{P}_B(1)$ constraint and
compares the posted bill with uniform pricing. Without bill protection, the
posted bill increases by 19.53\%; with protection, the posted bill does not
exceed the uniform-pricing level, while profit still improves by 21.29\%.
Thus, the 80.51\% profit gain in the main results should be interpreted as
a scheduling benefit under a fixed average price, not as a bill-neutrality
benefit.

\begin{table}[htbp]
\centering
\caption{Effect of the fixed-average-price constraint with and without bill
protection. The posted bill without protection exceeds the uniform baseline
by 19.53\%, while profit still improves by 21.29\% even under strict bill
protection ($\rho=1$).}
\label{tab:bill_protection}
\begin{tabular}{lrrrr}
\toprule
Policy & Profit & Posted bill & Bill ratio to uniform & Min QoS \\
\midrule
Uniform pricing & 3600.20 & 9748.05 & 1.0000 & 0.6632 \\
Fixed average price & 6498.68 & 11652.22 & 1.1953 & 0.9918 \\
Fixed avg.\ price + bill protection & 4366.70 & 9748.05 & 1.0000 & 0.9572 \\
\bottomrule
\end{tabular}
\end{table}

To provide an enumerable numerical reference, we reduce the number of
periods to $T=4$ and enumerate the 4495 grid points satisfying the mean
constraint at a step size of 0.05. The best grid profit is 2918.94, while
the multi-start SLSQP candidate profit is 2926.57 (0.261\% higher). This
check confirms that the candidate is at least no worse than the coarse grid.

\subsection{RQ2: When Does QoS Modeling Yield Incremental Value?}

\begin{table}[htbp]
\centering
\caption{Sensitivity to QoS degradation severity $\kappa$. The incremental
value of explicit QoS modeling is measured as the profit gain of QoS-aware
pricing over myopic pricing (parentheses). This gain increases from 0\% at
$\kappa=0$ to 12.53\% at $\kappa=30$.}
\label{tab:kappa}
\begin{tabular}{rrrr}
\toprule
$\kappa$ & Profit gain vs.\ uniform & Profit gain vs.\ myopic & Min QoS \\
\midrule
0  & 37.75\% & 0.00\% & 1.0000 \\
5  & 49.96\% & 0.99\% & 0.9871 \\
15 & 80.51\% & 5.22\% & 0.9918 \\
30 & 136.23\% & 12.53\% & 0.9952 \\
\bottomrule
\end{tabular}
\end{table}

Table~\ref{tab:kappa} reveals a boundary that must be stated explicitly:
even when $\kappa=0$, dynamic pricing yields a 37.75\% gain over uniform
pricing. This portion arises from heterogeneity in time-slot preferences,
price sensitivity, and rigid-user willingness to pay, not from QoS
degradation. The incremental value of QoS modeling should be measured by
the difference between QoS-aware and myopic dynamic pricing.

\subsection{RQ3: Are Conclusions Robust to Key Assumptions?}

\begin{table}[htbp]
\centering
\caption{Sensitivity of the QoS-aware profit increment to key assumptions.
The advantage over myopic dynamic pricing is largest under tight capacity
($G=750$, +293.78\%) and disappears under abundant capacity ($G=1200$, 0\%).}
\label{tab:ablation}
\begin{tabular}{llr}
\toprule
Experiment & Parameter value & Profit increment \\
\midrule
Fixed average price & $0.60 / 0.80 / 1.00 / 1.20$ & $9.90 / 5.22 / 2.46 / 2.74$\% \\
Capacity & $750 / 950 / 1200$ & $293.78 / 5.22 / 0.00$\% \\
Market expansion coefficient & $0 / 2.5 / 6.0$ & $4.88 / 5.22 / 5.71$\% \\
QoS feedback weight & $0 / 1 / 2$ & $5.22 / 4.99 / 4.80$\% \\
\bottomrule
\end{tabular}
\end{table}

When $G=750$, uniform pricing yields a negative profit, so percentage gains
over uniform pricing would be misleading. Table~\ref{tab:ablation} instead
compares QoS-aware pricing with myopic dynamic pricing. When $G=1200$, the
two policies achieve nearly identical profit, indicating that the incremental
value of explicit QoS modeling vanishes when capacity is abundant.

\subsection{Rule Baselines, Billing Ablation, and Scalability}

To verify that the QoS-aware multi-start search outperforms more than just
a weak baseline, we add peak--valley rules and queue-aware rules. Both rule
types adjust per-period prices based on the rigid load profile, then project
back onto the mean constraint via~\eqref{eq:projection}.
Table~\ref{tab:supplemental_policy} reports the results.

\begin{table}[htbp]
\centering
\caption{Rule-based pricing baselines and billing mechanism ablation. The
QoS-aware multi-start search outperforms both heuristic rules, and the
choice of billing rule (full billing vs.\ SLA penalty) affects the optimal
load distribution.}
\label{tab:supplemental_policy}
\small
\begin{tabular}{lrrr}
\toprule
Policy or billing mechanism & Profit & Welfare & Min QoS \\
\midrule
Uniform pricing & 3600.20 & 14307.89 & 0.6632 \\
Peak--valley rule & 5633.67 & 15331.25 & 0.8533 \\
Queue-aware rule & 6489.51 & 15697.14 & 0.9897 \\
QoS-aware multi-start search & \textbf{6498.68} & \textbf{15724.61} & \textbf{0.9918} \\
\midrule
Full billing & 6575.24 & 15815.07 & 0.9423 \\
SLA penalty billing & 6545.08 & 15786.57 & 0.9650 \\
\bottomrule
\end{tabular}
\end{table}

\subsection{Single-GPU System Measurements}
\label{sec:system_measurement}

Supplementary single-GPU load experiments use Ollama 0.24.0 on RTX 4090
with a 9B-parameter Q4\_K\_M quantized GGUF model.

\begin{table}[htbp]
\centering
\caption{Single-GPU Ollama load experiment on RTX 4090 with a 9B-parameter quantized model. Throughput saturates beyond two concurrent requests while TTFT increases steadily, indicating queueing-dominated congestion.}
\label{tab:system_measurement}
\small
\begin{tabular}{rrrrrr}
\toprule
Concurrency & Completion rate & Mean TTFT (s) & P95 TTFT (s) & TPOT (ms) &
Throughput (token/s) \\
\midrule
1 & 100\% & 0.155 & 0.167 & 9.456 & 91.46 \\
2 & 100\% & 1.232 & 1.421 & 9.413 & 99.28 \\
4 & 100\% & 3.296 & 4.032 & 9.426 & 98.75 \\
8 & 100\% & 6.136 & 9.175 & 9.464 & 98.63 \\
\bottomrule
\end{tabular}
\end{table}

Table~\ref{tab:system_measurement} shows that total throughput increases by
8.55\% when concurrency rises from 1 to 2, then largely saturates; mean
TTFT and P95 TTFT increase steadily with concurrency. The primary congestion
manifestation in this backend is queueing growth rather than TPOT
degradation.

\subsection{vLLM Primary Backend Overload Experiments}
\label{sec:vllm_measurement}

Repeated-overload scans using vLLM 0.22.0 with Qwen2.5-0.5B-Instruct on
RTX 4090.

\begin{table}[htbp]
\centering
\caption{vLLM repeated-overload scan on RTX 4090 with Qwen2.5-0.5B. The
0.5~s SLA attainment rate drops from 100.0\% at 64 concurrency to 50.0\% at
512 concurrency. Parentheses show 95\% confidence half-widths over 5
repetitions of 512 requests each.}
\label{tab:vllm_measurement}
\small
\begin{tabular}{rrrrr}
\toprule
Concurrency & Mean TTFT (s) & P95 TTFT (s) & Throughput (token/s) &
0.5~s SLA rate \\
\midrule
64  & 0.049 (0.001) & 0.073 (0.003) & 16125.98 (246.35) & 100.0 (0.0)\% \\
128 & 0.084 (0.004) & 0.139 (0.006) & 19854.91 (568.53) & 100.0 (0.0)\% \\
224 & 0.122 (0.005) & 0.225 (0.029) & 20128.54 (736.48) & 100.0 (0.0)\% \\
384 & 0.544 (0.023) & 1.617 (0.077) & 18983.53 (453.70) & 63.5 (0.7)\% \\
512 & 0.632 (0.038) & 1.588 (0.063) & 19266.09 (291.99) & 50.0 (0.0)\% \\
\bottomrule
\end{tabular}
\end{table}

\begin{figure}[htbp]
\centering
\includegraphics[width=\textwidth]{figures/vllm_reliability_curve.pdf}
\caption{vLLM overload curves for Qwen2.5-0.5B on RTX 4090. Left: mean and
P95 TTFT versus concurrency; right: throughput and 0.5~s SLA attainment
rate. The SLA rate falls sharply between 224 and 384 concurrent requests,
confirming a measurable capacity boundary. Error bars show 95\% confidence
half-widths.}
\label{fig:vllm_overload}
\end{figure}

\subsection{Empirical QoS Heuristic Mapping and Robustness Check}

Using TTFT SLA attainment rate as the empirical QoS observation, we
normalize the last concurrency level with an SLA rate of at least 99\% as
$u=1$, then construct heuristic scenario mappings via~\eqref{eq:qos}.

\begin{table}[htbp]
\centering
\caption{Empirical QoS heuristic mapping from vLLM repeated-overload scans.
The normalized healthy capacity is 224 concurrent requests for both models.
The 0.5B model shows higher fitting error (RMSE 0.0683) due to the 384-concurrency transition point.}
\label{tab:qos_fit}
\begin{tabular}{lrrr}
\toprule
Model & Normalized threshold & Degradation severity & RMSE \\
\midrule
Qwen2.5-0.5B-Instruct & 1.000 & 0.517 & 0.0683 \\
Qwen2.5-3B-Instruct & 1.058 & 0.942 & $6.33\times10^{-10}$ \\
\bottomrule
\end{tabular}
\end{table}

Substituting both sets of heuristic parameters into the economic model and
re-optimizing prices yields a profit of 6649.24 and a minimum QoS of 1.0000
in both cases, compared to 6498.68 for the baseline synthetic proxy.

\section{Discussion}
\label{sec:discussion}

The current version primarily uses algorithmic simulation, supplemented by
single-GPU backend measurements using vLLM and Ollama to validate the
congestion phenomenon. These measurements cannot substitute for system
experiments across a broader range of model sizes, multiple backends, and
multi-GPU configurations. The electricity dynamic pricing literature
provides a methodological reference, but inference serving does not involve
grid power flow, nodal balance, or wholesale market clearing constraints.

Key limitations include: (1) the QoS function remains a parametric proxy;
(2) the welfare metric depends on synthetic parameters; (3) multi-start
SLSQP does not provide a global optimality certificate for the non-convex
problem; (4) the aggregate logit response is not a proof of pure-strategy
Nash equilibrium; (5) multi-platform competition is excluded; (6) the
demand curve is derived from synthetic parameters; (7) a fixed average
price is not equivalent to bill neutrality; (8) price-sensitivity analysis
does not demonstrate distributional fairness.

\section{Conclusion}

We adapt the intertemporal scarcity-pricing mechanism and leader--follower
hierarchical structure from electricity demand response to token services.
The platform selects a price vector under a fixed-average-price constraint,
and elastic users migrate demand in response to prices and time-slot
preferences. On the theoretical side, we prove the existence of a nonempty
feasible price region and an optimal platform solution; on the algorithmic
side, we implement an exact bounded-mean projection, multi-start
optimization, and traceable experimental artifacts. Synthetic experiments
answer three research questions: constrained dynamic pricing can shift
demand and alleviate congestion; the profit advantage of explicit QoS
modeling over myopic pricing increases with degradation severity and
disappears when capacity is abundant; and this incremental value maintains
interpretable boundaries across multiple sensitivity experiments. Under
the additional bill-protection constraint, profit still improves by 21.29\%
over uniform pricing. These conclusions are confined to synthetic demand
and a parametric QoS proxy.

\section{Code and Artifact Availability}

All code, experimental configurations, and analysis scripts required to
reproduce the results are included in the supplementary materials. The
official public repository URL will be made available upon publication.

\appendix
\section{Baseline Parameters}

\begin{table}[H]
\centering
\caption{Synthetic economic model baseline parameters.}
\label{tab:baseline_parameters}
\begin{tabular}{lll}
\toprule
Parameter & Baseline value & Description \\
\midrule
$T,h$ & $8,3$ & Number of periods and hours per period \\
$G,w$ & $950,0.40$ & Capacity and unit serving cost \\
$\pcap,\underline p,\overline p$ & $0.80,0.45,2.10$ & Mean price and price bounds \\
$\overline u,\kappa$ & $0.82,15$ & QoS threshold and degradation severity \\
$\alpha,\beta$ & $3.5,5.0$ & Elastic and rigid price sensitivity \\
$v_r,\overline{\F}$ & $1.80,400$ & Rigid willingness-to-pay ceiling and elastic market size \\
$c_s,\eta$ & $0.20,2.50$ & Switching cost and market expansion coefficient \\
$c_r,c_q$ & $0.60,0.35$ & Rigid churn cost and QoS degradation cost \\
\bottomrule
\end{tabular}
\end{table}

The rigid load vector and time-slot preference vector are
\[
(200,150,300,600,800,900,700,250),\qquad
(0.3,0.2,0.5,0.8,1.0,1.0,0.9,0.4),
\]
respectively. The carbon emission rate and unit carbon cost are 0.20 and
0.05. The bill-protection scenario uses $\rho=1$.

\section{Complete Sensitivity Range}

\begin{table}[H]
\centering
\caption{Phase one complete sensitivity range.}
\label{tab:full_sensitivity}
\begin{tabular}{lll}
\toprule
Experiment & Parameter range & Primary purpose \\
\midrule
Fixed average price & $0.60,0.80,1.00,1.20$ & Examine price constraint boundary \\
QoS degradation severity & $0,5,15,30$ & Separate dynamic pricing gains from QoS incremental value \\
QoS threshold & $0.75,0.82,0.90$ & Examine congestion trigger location \\
Capacity & $750,950,1200$ & Examine tight vs.\ abundant capacity \\
Price sensitivity & $1.0,3.5,8.0$ & Examine elastic user response \\
Switching cost & $0.05,0.20,0.80$ & Examine intertemporal migration resistance \\
Market expansion coefficient & $0,2.5,6.0$ & Rule out market expansion as the sole source \\
Churn cost & $0,0.60,1.20$ & Examine long-term cost assumptions \\
QoS feedback weight & $0,1,2$ & Examine demand feedback variant \\
\bottomrule
\end{tabular}
\end{table}

\small
\bibliographystyle{unsrtnat}
\bibliography{verified_refs}

\end{document}
